\section{Moment batches and their full gauge}\label{gauge:section}

We now turn the positive finite moments into a deterministic gauge.
Each finite batch must balance two requirements: moderately large masses
should provide many chances for a successful test, whereas exceeding a
larger multiple of the threshold should have summable probability.
Deterministic repetitions of dyadic thresholds achieve both.

Write \(\PP^*_{\mathrm{mix}}\) for the law of the disk model in
Equation~\eqref{coef:disk-model}, with initial angle distributed according
to \(\pi_a\), and write \(\EE^*_{\mathrm{mix}}\) for its expectation.
We use the increasing masses \(L_\ell\uparrow L\) and their moments
\(m_n(\ell),m_n\) from Equation~\eqref{coef:moments}. All the numbers below
are deterministic: the moments are specified by the explicit integrals
in Equation~\eqref{coef:formula}.

\subsection{The finite lists and their radii}

For integers \(n\ge1\) and \(k\in\mathbb Z\), put
\begin{equation}\label{gauge:coefficients}
\ell_n=1+\sum_{j=1}^\infty
\ind_{\{m_n(j)\le 3m_n/4\}},
\qquad A_k=2^k,
\qquad T_n=\frac{4m_{n+1}}{m_n}.
\end{equation}
The sum is finite, and \(m_n(\ell_n)>3m_n/4\), since
\(m_n(j)\uparrow m_n\in(0,\infty)\). Form a list, in increasing
order of \(k\), containing \(A_k\) with multiplicity
\begin{equation}\label{gauge:batch}
N_{nk}=\ind_{\{2^k\le T_n\}}
\left\lfloor\frac{4^n2^{kn}}{m_n}\right\rfloor.
\end{equation}
Append the number \(n\) once, as the final \emph{cleanup entry},
and denote the resulting entries by \(b_{ni}\), \(1\le i\le I_n\).
This list is finite. Indeed, \(N_{nk}>0\) implies
\[
\frac{m_n^{1/n}}4\le A_k\le T_n,
\]
and Jensen's Inequality gives \(m_n^{1/n}\ge m_1>0\).
Thus \(c_0:=\min\{1,m_1/4\}\) is a positive lower bound for every
entry \(b_{ni}\).

Assign the radii
\begin{equation}\label{gauge:radii}
r_{ni}=\exp\{-n^2-2a(i-1)(\ell_n+n)\},
\qquad 1\le i\le I_n,
\end{equation}
and define
\begin{equation}\label{gauge:definition}
\boxed{\quad
h_\kappa(r)=
\inf_{\substack{n\ge1\\1\le i\le I_n}}
b_{ni}r_{ni}^{\,d}\max\{1,(r/r_{ni})^2\}
\quad(r>0),\qquad h_\kappa(0)=0.
\quad}
\end{equation}
The infimum includes every entry of every batch.

The quadratic envelope has a planar covering antecedent in
Peres--Solomyak \cite[Lemma~1.2 and Section~5]{PeresSolomyak2005}.
Indeed, for $0<s\le e^{-1}$ put
$\phi(s)=\inf_{r_{ni}\ge s}b_{ni}r_{ni}^{d}$. Exchanging infima gives
\[
 h_\kappa(s)=\inf_{0<r\le s}(s/r)^2\phi(r),
\]
which is their quadratic regularization of this step profile.
The batch data, nondegeneracy of the envelope, and control after finite
exclusions are established here; in particular, positivity is proved
directly rather than inferred for an arbitrary input profile.

\subsection{Admissibility and restoration of early entries}

The envelope must remain positive at every positive diameter and tend
to zero at the origin. Its behavior after a finite set of entries is
excluded will also be needed in the lower-cover proof.

\begin{proposition}[Gauge properties and finite exclusions]\label{gauge:admissible}
The function \(h_\kappa\) is a Hausdorff gauge and satisfies
\begin{equation}\label{gauge:doubling}
h_\kappa(Cr)\le C^2h_\kappa(r)
\qquad(C\ge1,\ r\ge0).
\end{equation}
The radii from distinct batches may interleave. Set
\[
H_{ni}(r)=b_{ni}r_{ni}^{\,d}\max\{1,(r/r_{ni})^2\}.
\]
If \(0\le F(r)\le M<\infty\) and
\(F(r)\le C H_{ni}(r)\) for every \(n\ge n_0\), every
\(1\le i\le I_n\), and every \(r>0\), then
\(F(r)\le C' h_\kappa(r)\) for all \(r>0\), for some finite
\(C'\).
\end{proposition}

\begin{proof}
There are only finitely many entries with \(r_{ni}\ge\varepsilon\)
for each \(\varepsilon>0\), since \(r_{ni}\le e^{-n^2}\) and
every batch is finite. For a fixed \(r>0\), entries with
\(r_{ni}<r\) satisfy
\[
H_{ni}(r)\ge c_0r^2r_{ni}^{d-2}\longrightarrow\infty
\qquad\text{as }r_{ni}\longrightarrow0,
\]
because \(d<2\). Comparing with any one entry reduces the infimum
to finitely many positive terms. Hence \(0<h_\kappa(r)<\infty\).

Each \(H_{ni}\) is nondecreasing and satisfies
\(H_{ni}(Cr)\le C^2H_{ni}(r)\). Taking infima proves
monotonicity and Equation~\eqref{gauge:doubling}. In particular,
for \(0<s<r\),
\[
(s/r)^2h_\kappa(r)\le h_\kappa(s)\le h_\kappa(r),
\]
which gives continuity on \((0,\infty)\). At the cleanup radius
\(\rho_n=r_{nI_n}\),
\[
h_\kappa(\rho_n)\le n\rho_n^d\le ne^{-dn^2}\longrightarrow0.
\]
Since \(\rho_n\to0\), monotonicity gives continuity at zero.

For the last assertion the discarded entries, those with \(n<n_0\),
are a finite set, and \(H_{ni}(r)\ge b_{ni}r_{ni}^d>0\).
Thus one may take
\[
C'=\max\left\{C,
\max_{\substack{1\le n<n_0\\1\le i\le I_n}}
\frac{M}{b_{ni}r_{ni}^d}\right\},
\]
omitting the inner maximum if its index set is empty.
Then \(F(r)\le C'H_{ni}(r)\) holds for every entry, and taking
the full infimum proves the claim. No ordering between batches
is used.
\end{proof}

\subsection{The opposing probability estimates}

The multiplicities supply the two bounds advertised at the start of
this section. These are estimates under the one mixed disk law; the
conditional tests along the curve are constructed in
Section~\ref{hd:cover-section}.

\begin{lemma}[Batch probability estimates]\label{gauge:multiplicity}
The non-cleanup multiplicities satisfy, for every \(n\ge1\),
\begin{equation}\label{gauge:upper-tail}
\sum_{k\in\mathbb Z}N_{nk}\,
\PP^*_{\mathrm{mix}}\{L\ge16A_k\}
\le\frac{4^n16^{-n}}{1-2^{-n}}.
\end{equation}
For \(n\ge3\), they also satisfy
\begin{equation}\label{gauge:success}
\sum_{k\in\mathbb Z}N_{nk}\,
\PP^*_{\mathrm{mix}}\{L_{\ell_n}\ge A_k\}
\ge 2^{n-2}-\frac{1}{2(1-2^{-n})}
\ge\frac{2^n}{8}.
\end{equation}
Including the cleanup entries,
\begin{equation}\label{gauge:tail-summability}
\sum_{n\ge1}\sum_{i=1}^{I_n}
\PP^*_{\mathrm{mix}}\{L\ge16b_{ni}\}<\infty.
\end{equation}
\end{lemma}

\begin{proof}
For \(x\ge0\), the geometric sum gives
\[
\sum_{k:\,2^k\le x}2^{kn}
\le\frac{x^n}{1-2^{-n}}.
\]
Discarding the cutoff and the floors in Equation~\eqref{gauge:batch},
then applying this inequality with \(x=L/16\) and taking
expectations, proves Equation~\eqref{gauge:upper-tail}.

For the lower bound put \(Y=L_{\ell_n}\). Since \(Y\le L\),
\[
\EE^*_{\mathrm{mix}}[Y^n;Y>T_n]
\le\frac{\EE^*_{\mathrm{mix}}[Y^{n+1}]}{T_n}
\le\frac{m_n}{4},
\qquad
\EE^*_{\mathrm{mix}}[Y^n;Y\le T_n]\ge\frac{m_n}{2}.
\]
Write \(w_k=4^nA_k^n/m_n\). Whenever \(0<Y\le T_n\),
some dyadic level satisfies \(Y/2<A_k\le Y\); it is allowed
by the cutoff, and its weight is at least \(2^nY^n/m_n\).
Consequently
\[
\EE^*_{\mathrm{mix}}\left[
\sum_{k:\,A_k\le T_n}w_k\ind_{\{Y\ge A_k\}}\right]
\ge2^{n-1}.
\]
The weights with \(w_k<1\) have total at most
\((1-2^{-n})^{-1}\), by another geometric sum. For the other
weights, \(\lfloor w_k\rfloor\ge w_k/2\). Subtracting the former
weights and halving therefore gives
\[
\sum_kN_{nk}\PP^*_{\mathrm{mix}}\{Y\ge A_k\}
\ge2^{n-2}-\frac{1}{2(1-2^{-n})}.
\]
For \(n\ge3\), the subtracted term is at most \(4/7\le2^{n-3}\),
which proves Equation~\eqref{gauge:success}.
Finally, Tonelli's Theorem gives
\[
\sum_{n\ge1}\PP^*_{\mathrm{mix}}\{L\ge16n\}
=\EE^*_{\mathrm{mix}}\!\left[\left\lfloor\frac L{16}\right\rfloor\right]
\le\frac{m_1}{16}.
\]
Together with the summable bounds in Equation~\eqref{gauge:upper-tail},
this proves Equation~\eqref{gauge:tail-summability}.
\end{proof}
